% =============================================================================
%  style.tex — the paper's design system.
%
%  Single-column modern-preprint layout: XCharter for text, Fira Sans for
%  headings and captions, Inconsolata for code. Colours are the same three
%  categorical slots the figures use (paper/figures/style.py), so a reader
%  never has to relearn what blue and orange mean between a figure and the
%  prose around it.
% =============================================================================

% ---------------------------------------------------------------- geometry --
\usepackage[
  letterpaper,
  top=1.05in, bottom=1.05in, left=1.10in, right=1.10in,
  footskip=0.45in, headsep=0.22in
]{geometry}

% ------------------------------------------------------------------- fonts --
\usepackage[T1]{fontenc}
\usepackage{XCharter}
\usepackage[scaled=0.93]{FiraSans}
\usepackage[varqu,varl,scaled=0.92]{inconsolata}
\usepackage[charter,vvarbb,scaled=1.04]{newtxmath}
\usepackage{microtype}

\linespread{1.045}
\setlength{\emergencystretch}{1.6em}   % long \texttt{} runs must not overflow

% ------------------------------------------------------------------ colours --
\usepackage[table,dvipsnames]{xcolor}

\definecolor{arblue}{HTML}{2A78D6}    % autoregressive  (categorical slot 1)
\definecolor{dgorange}{HTML}{EB6834}  % block diffusion (categorical slot 2)
\definecolor{llaqua}{HTML}{1BAF7A}    % LLaDA           (categorical slot 3)

\definecolor{ink}{HTML}{0B0B0B}
\definecolor{ink2}{HTML}{52514E}
\definecolor{muted}{HTML}{898781}
\definecolor{rule}{HTML}{C3C2B7}
\definecolor{hairline}{HTML}{E8E7E1}

\definecolor{washblue}{HTML}{EAF1FB}
\definecolor{washorange}{HTML}{FDEEE7}
\definecolor{washneutral}{HTML}{F4F3F0}

\definecolor{deepink}{HTML}{1B3A5C}    % heading ink — a calm dark blue

% ------------------------------------------------------------------ tables --
\usepackage{booktabs}
\usepackage{array}
\usepackage{multirow}
\usepackage{tabularx}
\renewcommand{\arraystretch}{1.18}
\setlength{\tabcolsep}{6pt}
\newcolumntype{R}{>{\raggedleft\arraybackslash}X}
\newcolumntype{L}{>{\raggedright\arraybackslash}X}

% Table body in the sans face at a slightly smaller size — numbers read
% better and the table separates cleanly from the surrounding text.
\newcommand{\tabfont}{\sffamily\footnotesize}
\newcommand{\thead}[1]{\textcolor{ink2}{\footnotesize\textbf{#1}}}

% --------------------------------------------------------------- headings ---
\usepackage{titlesec}

\titleformat{\section}
  {\normalfont\sffamily\large\bfseries\color{deepink}}
  {\textcolor{arblue}{\thesection}}{0.7em}{}
\titlespacing*{\section}{0pt}{1.55\baselineskip}{0.55\baselineskip}

\titleformat{\subsection}
  {\normalfont\sffamily\normalsize\bfseries\color{deepink}}
  {\textcolor{muted}{\thesubsection}}{0.6em}{}
\titlespacing*{\subsection}{0pt}{1.15\baselineskip}{0.35\baselineskip}

\titleformat{\subsubsection}
  {\normalfont\sffamily\small\bfseries\color{ink2}}
  {}{0em}{}
\titlespacing*{\subsubsection}{0pt}{0.9\baselineskip}{0.25\baselineskip}

\titleformat{\paragraph}[runin]
  {\normalfont\sffamily\small\bfseries\color{ink}}
  {}{0em}{}[.]
\titlespacing*{\paragraph}{0pt}{0.7\baselineskip}{0.55em}

% --------------------------------------------------------------- captions ---
\usepackage[font=small,labelfont={sf,bf,color=deepink},
            textfont={color=ink2},labelsep=period,
            skip=7pt,justification=justified]{caption}
\usepackage{subcaption}

% ------------------------------------------------------------------- boxes --
\usepackage[most]{tcolorbox}

% A key-result callout: tinted panel, coloured left rule, sans heading.
\newtcolorbox{keybox}[1][]{
  enhanced, breakable,
  colback=washneutral, colframe=washneutral,
  boxrule=0pt, leftrule=2.2pt,
  borderline west={2.2pt}{0pt}{arblue},
  arc=0pt, outer arc=0pt,
  left=9pt, right=9pt, top=7pt, bottom=7pt,
  before skip=9pt, after skip=9pt,
  fontupper=\small,
  #1
}

% A "what this section establishes" strip, used once per major section.
\newtcolorbox{takeaway}[1][]{
  enhanced, breakable,
  colback=washblue!55!white, colframe=washblue!55!white,
  boxrule=0pt, arc=1.5pt, outer arc=1.5pt,
  left=9pt, right=9pt, top=6pt, bottom=6pt,
  before skip=9pt, after skip=9pt,
  fontupper=\small,
  #1
}

% A caution/limitation strip — for scope conditions the reader must carry.
\newtcolorbox{caution}[1][]{
  enhanced, breakable,
  colback=washorange!60!white, colframe=washorange!60!white,
  boxrule=0pt, leftrule=2.2pt,
  borderline west={2.2pt}{0pt}{dgorange},
  arc=0pt, outer arc=0pt,
  left=9pt, right=9pt, top=6pt, bottom=6pt,
  before skip=9pt, after skip=9pt,
  fontupper=\small,
  #1
}

\newcommand{\boxhead}[1]{%
  \textcolor{deepink}{\sffamily\bfseries\footnotesize\MakeUppercase{#1}}\par
  \vspace{2pt}}

% ------------------------------------------------------------------- lists --
\usepackage{enumitem}
\setlist[itemize]{leftmargin=1.15em,itemsep=2.5pt,topsep=4pt,parsep=0pt}
\setlist[enumerate]{leftmargin=1.5em,itemsep=2.5pt,topsep=4pt,parsep=0pt}

% -------------------------------------------------------------------- misc --
\usepackage{amsmath}
\let\Bbbk\relax          % newtxmath already defines it
\usepackage{amssymb}
\usepackage{siunitx}
\sisetup{group-separator={,},group-minimum-digits=4,detect-all}
\usepackage{graphicx}
\usepackage{algorithm}
\usepackage{algpseudocode}
\usepackage{float}
\usepackage[hidelinks]{hyperref}
\hypersetup{
  colorlinks=true,
  linkcolor=arblue, citecolor=arblue, urlcolor=arblue,
  pdfborder={0 0 0}
}

% Model shorthands — one spelling per model, everywhere.
\newcommand{\ARthree}{\textsc{ar-3b}}
\newcommand{\ARoneb}{\textsc{ar-1.5b}}
\newcommand{\DG}{\textsc{dg}}
\newcommand{\LLaDA}{\textsc{llada-8b}}

% Coloured inline model tags, matching the figures.
\newcommand{\artag}{\textcolor{arblue}{\sffamily\bfseries AR}}
\newcommand{\dgtag}{\textcolor{dgorange}{\sffamily\bfseries DG}}

% A quiet inline source pointer.
\newcommand{\artifact}[1]{\textcolor{muted}{\ttfamily\scriptsize #1}}

% ---------------------------------------------------------------- title page --
\usepackage{titling}
\usepackage{fancyhdr}

\pagestyle{fancy}
\fancyhf{}
\renewcommand{\headrulewidth}{0pt}
\fancyfoot[C]{\textcolor{muted}{\sffamily\footnotesize\thepage}}

% \papertitle{title}{subtitle}{author}{affiliation}{version line}
\newcommand{\papertitle}[5]{%
  \thispagestyle{empty}
  \begingroup
  \raggedright\parindent=0pt
  {\sffamily\bfseries\fontsize{21}{25}\selectfont\color{deepink} #1\par}
  \vspace{5pt}
  {\sffamily\fontsize{13}{16}\selectfont\color{ink2} #2\par}
  \vspace{8pt}
  {\color{arblue}\rule{\linewidth}{2pt}\par}
  \vspace{9pt}
  {\sffamily\large\color{ink} #3\par}
  \vspace{2pt}
  {\sffamily\small\color{ink2} #4\par}
  \vspace{3pt}
  {\sffamily\footnotesize\color{muted} #5\par}
  \endgroup
  \vspace{16pt}
}
